\chapter{A Question's Journey: End-to-End Execution Trace}
\label{ch:trace}

This chapter follows one real question through every function of the shipped \code{guarded} profile, using measurements captured during this analysis rather than illustrations. The run happened to include a failure, which makes it more instructive than a clean run would have been.

\begin{measured}[how the trace was captured]
A script on the development machine built \code{Pipeline(Settings(), load_pipeline_config("guarded"))} against the working index, ran one retrieval to load the models, then timed \code{Pipeline.answer("What happens during anaphase of mitosis?")}. It wrapped the Gemini client's \code{generate_content} to record each call's duration and token usage, and wrapped the grader to keep its claims. No project code was changed. The question is one of the demo's examples.
\end{measured}

\section{The result at a glance}

\begin{center}
\begin{tabularx}{\textwidth}{@{}>{\raggedright\arraybackslash}p{5.2cm}L@{}}
\toprule
Model loading (first retrieval) & 23,723\,ms, before the question was asked\\
Wall-clock time for \code{answer()} & 19,184\,ms\\
Reported \code{trace.total_ms} & 37,431\,ms --- the double count of Section~\ref{sec:latency-bug}\\
Stages recorded & \code{guard_input} 0.1, \code{retrieve} 1,569.7, \code{rerank} 14,870.7, \code{guard_context} 24.2, \code{generate} 1,782.0, \code{corrective} 19,183.8 (ms)\\
Model calls & Generation: 2,367 input and 156 output tokens, recorded. Grading: \textbf{failed with 503 UNAVAILABLE}\\
Decision & ACCEPT, on rule 1 (grading not measured)\\
Answer & Four sentences citing passages 2, 3, 4 and 5; \code{grounded=True}\\
\bottomrule
\end{tabularx}
\end{center}

\begin{figure}[htbp]
\centering
\resizebox{\textwidth}{!}{%
\begin{tikzpicture}[font=\sffamily\scriptsize, x=0.42cm, y=0.62cm]
\draw[-Stealth] (0,0) -- (39.5,0) node[right] {s};
\foreach \t in {0,5,10,15,20,25,30,35} { \draw (\t,0.12) -- (\t,-0.12) node[below] {\t}; }
\fill[blue!22] (0,5) rectangle (19.18,5.7); \node[anchor=west] at (19.5,5.35) {\code{corrective} 19,184 ms};
\fill[orange!55] (0,3.9) rectangle (1.57,4.6); \node[anchor=south] at (0.8,4.6) {retrieve};
\fill[orange!35] (1.57,3.9) rectangle (16.44,4.6); \node at (9.0,4.25) {rerank 14,871 ms};
\fill[orange!55] (16.44,3.9) rectangle (18.25,4.6); \node[anchor=south] at (17.3,4.6) {generate};
\fill[red!30] (18.25,2.8) rectangle (19.18,3.5); \node[anchor=west] at (19.5,3.15) {grading call: 503 after about 0.9 s};
\draw[very thick, green!45!black] (0,1.9) -- (19.18,1.9); \node[anchor=west] at (19.5,1.9) {wall clock 19,184 ms};
\draw[very thick, red!70!black] (0,1.0) -- (37.43,1.0); \node[anchor=south east] at (37.43,1.05) {reported \code{total_ms} 37,431 ms};
\end{tikzpicture}}
\caption{The traced query on a time axis. The context guard (24\,ms) and input guard (0.1\,ms) are too short to draw.}
\label{fig:trace-timeline}
\end{figure}

\section{Step by step}

\subsection{Step 0 --- before the question: loading models}

Neither model is loaded when the pipeline is built. The first call to \code{embed_texts} creates the fastembed \code{TextEmbedding} and the first call to \code{rerank} creates the \code{TextCrossEncoder}; both are then cached for the life of the process by \code{lru_cache} \ev{src/vidyarag/llm/provider.py:33-47}, \ev{src/vidyarag/retrieve/rerank.py:31-41}. The warm-up retrieval that did this took 23.7 seconds (a later benchmark, with the same model files already on disk, loaded both models in about 5 seconds). A user whose question is the first after a restart waits for this too.

\subsection{Step 1 --- \code{Pipeline.answer} creates the trace}

\code{trace = QueryTrace(query=question, profile="guarded")} \ev{src/vidyarag/pipeline.py:153}. Every later step writes into this object.

\subsection{Step 2 --- the input guard: 0.1\,ms}

\code{check_user_input} is true, so \code{screen_input} runs the three input pattern families over the question inside a \code{guard_input} stage \ev{src/vidyarag/pipeline.py:159-171}. Nothing matched, so \code{trace.guard_input} stays empty and execution continues.

\subsection{Step 3 --- entering the self-check loop}

\code{corrective.enabled} is true, so \code{_answer_corrective} builds a \code{CorrectivePolicy} with an accept threshold of 0.8, an abstain threshold of 0.5 and a budget of two attempts, defines the \code{retrieve} and \code{generate} closures, opens the \code{corrective} stage and calls \code{run_corrective_loop} \ev{src/vidyarag/pipeline.py:203-238}. Attempt 1 begins with \code{query} equal to the question.

\subsection{Step 4 --- dense retrieval: 1,570\,ms}

\code{Pipeline.retrieve} opens the \code{retrieve} stage and calls \code{retrieve_dense}: the question is embedded into 768 numbers by the cached BGE model, and Qdrant returns the 20 most similar passages with their payloads \ev{src/vidyarag/retrieve/dense.py:105-124}. Their ids are saved in \code{trace.retrieved_chunk_ids}. The first ten, in order:

\begin{center}\small
\begin{tabular}{@{}rl@{\qquad}rl@{}}
\toprule
\textbf{Rank} & \textbf{Chunk id} & \textbf{Rank} & \textbf{Chunk id}\\
\midrule
1 & \code{anatomy-and-physiology-p0127-03} & 6 & \code{anatomy-and-physiology-p0128-04}\\
2 & \code{biology-p0288-03} & 7 & \code{biology-p0308-02}\\
3 & \code{biology-p0285-00} & 8 & \code{biology-p0302-00}\\
4 & \code{biology-p0286-02} & 9 & \code{biology-p0286-01}\\
5 & \code{biology-p0303-02} & 10 & \code{anatomy-and-physiology-p0126-01}\\
\bottomrule
\end{tabular}
\end{center}

\subsection{Step 5 --- reranking: 14,871\,ms}

\code{rerank} scores all 20 (question, passage) pairs with the cross-encoder and sorts them; \code{rank_movement} summarises the change as \code{{reordered: 17, promoted_into_context: 2, top1_changed: false}}, which is stored in \code{trace.rerank} \ev{src/vidyarag/pipeline.py:120-130}. The five passages that will reach the prompt:

\begin{center}\small
\begin{tabularx}{\textwidth}{@{}c>{\raggedright\arraybackslash}p{4.9cm}Lrrc@{}}
\toprule
\textbf{Final} & \textbf{Chunk id} & \textbf{Citation} & \textbf{Cross-encoder} & \textbf{Dense} & \textbf{Dense rank}\\
\midrule
1 & \code{anatomy-and-physiology-p0127-03} & Anatomy and Physiology, 3.5. Cell Growth and Division*, p.117 & 3.698 & 0.797 & 1\\
2 & \code{biology-p0286-02} & Biology, 10.2. The Cell Cycle*, p.274 & 3.496 & 0.753 & 4\\
3 & \code{biology-p0288-03} & Biology, 10.2. The Cell Cycle*, p.276 & 3.335 & 0.770 & 2\\
4 & \code{biology-p0302-00} & Biology, Glossary, p.290 & 3.199 & 0.747 & 8\\
5 & \code{anatomy-and-physiology-p0128-04} & Anatomy and Physiology, 3.5. Cell Growth and Division*, p.118 & 1.726 & 0.752 & 6\\
\bottomrule
\end{tabularx}
\end{center}

Two passages were promoted into the context (dense ranks 8 and 6), and two left it: \code{biology-p0285-00} (dense rank 3) and \code{biology-p0303-02} (dense rank 5, now 8). The cross-encoder scores are unnormalised, so only their order matters. Reranking took 77\% of this query's wall-clock time. That share was unusually high: a steady-state benchmark on the same machine later measured 2.1 seconds for the same 20 passages (Chapter~\ref{ch:performance}), so the machine was evidently slower during this run --- a reminder that single CPU timings on a laptop are fragile.

\subsection{Step 6 --- the context guard: 24\,ms}

The first five passages are scanned for embedded directives and instruction overrides inside a \code{guard_context} stage \ev{src/vidyarag/pipeline.py:136-140}. Nothing was quarantined, so \code{trace.guard_context} stays empty and all five passages continue.

\subsection{Step 7 --- generation: 1,782\,ms}

The \code{generate} closure calls \code{generate_answer} with the \emph{original} question. \code{build_answer_prompt} produces a 745-character system instruction and a 10,609-character user message --- five passages of about 2,000 characters each, numbered and labelled:

\begin{lstlisting}[language=plain]
Passages:

[1] Anatomy and Physiology, 3.5. Cell Growth and Division*, p.117
Mitosis and Cytokinesis The mitotic phase of the cell typically takes between 1 and 2 hours.
During this phase, a cell undergoes two major processes. First, it completes mitosis, during
which the contents of the nucleus are ...
\end{lstlisting}

Inside the \code{generate} stage, \code{client.models.generate_content} sends it to \code{gemini-3.5-flash-lite} at temperature 0. \code{_record_usage} adds 2,367 input and 156 output tokens to the trace; \code{_extract_text} returns the answer \ev{src/vidyarag/generate/answer.py:103-128}:

\begin{quote}\small
During anaphase, the connection or cohesin proteins between the sister chromatids break down and degrade, allowing the sister chromatids to separate from each other at the centromere and form individual chromosomes once again [2, 3, 4, 5]. The microtubules shorten, and the chromosomes are pulled rapidly toward opposite poles or the centrosome to which their microtubule is attached by their kinetochores [2, 3, 5]. This process ensures that each end of the cell receives one partner from each pair of sister chromatids, providing identical genetic material to the two new daughter cells [5]. Additionally, the cell becomes visibly elongated into an oval shape as the polar microtubules slide against each other where they overlap at the metaphase plate [3].
\end{quote}

\code{resolve_citations} finds markers 2, 3, 4 and 5, all within the five supplied passages, so four citations are created and none are stripped. Passage 1 --- the top-ranked one --- is not cited.

\begin{measured}[the same question, asked twice]
An earlier timing run on the same question also sent 2,367 input tokens, because retrieval is deterministic and produced the same prompt, but received 130 output tokens instead of 156. Identical input at temperature 0 gave a different answer length.
\end{measured}

\subsection{Step 8 --- grading: the call failed}
\label{sec:trace-failopen}

The loop called \code{grade_answer} with the draft and the five passage texts. The call to \code{gemini-3.1-flash-lite} raised a server error: \code{503 UNAVAILABLE}, \enquote{This model is currently experiencing high demand. Spikes in demand are usually temporary. Please try again later.} \code{grade_answer} caught it and returned \code{Groundedness(claims=[], error="ServerError: 503 UNAVAILABLE. ...")} \ev{src/vidyarag/correct/grader.py:206-217}. The call lasted about 0.9 seconds: the corrective stage minus the four stages nested inside it.

\code{policy.decide} then applied its first rule --- not measured, so ACCEPT \ev{src/vidyarag/correct/policy.py:68-69} --- and the loop returned the draft as the answer. The trace records exactly what happened:

\begin{lstlisting}[language=json]
"corrective": {
  "attempts": 1, "abstained": false, "fired": false, "final_score": null,
  "trace": [{
    "query": "What happens during anaphase of mitosis?",
    "decision": "accept", "score": 0.0, "claims": 0,
    "supported": 0, "partial": 0, "unsupported": 0,
    "refuses": false, "measured": false,
    "error": "ServerError: 503 UNAVAILABLE. {'error': {'code': 503, 'message': 'This model is currently experiencing high demand. ..."
  }]
}
\end{lstlisting}

\paragraph{What grading would have said.} To see the check that did not run, the same answer and the same five passages were graded again with \code{grade_answer} a few minutes later. The first retry also returned 503 (after 1,077\,ms); the next succeeded in 4,569\,ms, using 2,497 input and 612 output tokens:

\begin{center}\small
\begin{tabularx}{\textwidth}{@{}c>{\raggedright\arraybackslash}p{7.4cm}L@{}}
\toprule
\textbf{Verdict} & \textbf{Claim extracted by the grader} & \textbf{Evidence it quoted (truncated)}\\
\midrule
supported & During anaphase, the connection or cohesin proteins between the sister chromatids break down and degrade. & \enquote{In anaphase, the connection between the sister chromatids breaks down, and the microtubules pull the chromosom\ldots}\\
supported & The sister chromatids separate from each other at the centromere and form individual chromosomes once again. & \enquote{During anaphase, the \enquote{upward phase,} the cohesin proteins degrade, and the sister chromatids separate at the c\ldots}\\
supported & The microtubules shorten, and the chromosomes are pulled rapidly toward opposite poles or the centrosome to which their microtubule is attached by their kinetochores. & \enquote{Each chromatid, now called a chromosome, is pulled rapidly toward the centrosome to which its microtubule is a\ldots}\\
supported & This process ensures that each end of the cell receives one partner from each pair of sister chromatids, providing identical genetic material to the two new daughter cells. & \enquote{Each end of the cell receives one partner from each pair of sister chromatids, ensuring that the two new daugh\ldots}\\
supported & The cell becomes visibly elongated into an oval shape as the polar microtubules slide against each other where they overlap at the metaphase plate. & \enquote{The cell becomes visibly elongated (oval shaped) as the polar microtubules slide against each other at the met\ldots}\\
\bottomrule
\end{tabularx}
\end{center}

Five claims, all supported: $g = 5/5 = 1.0$, \code{refuses=false}. Had grading succeeded, the policy would have accepted the same answer on rule 3. The outcome was right; the verification simply did not happen, and nothing outside the trace says so.

\subsection{Step 9 --- building the result}

The loop returned \code{LoopOutcome(answer=<draft>, abstained=False, attempts=[1 attempt])}. \code{_answer_corrective} copied the attempt count, the abstention flag and the loop record into the trace, and --- because the answer was not an abstention --- returned an \code{Answer} carrying the last draft's four citations, the five passages and \code{grounded=True} \ev{src/vidyarag/pipeline.py:240-267}.

\subsection{Step 10 --- what each interface would show}

\begin{description}[style=nextline]
  \item[HTTP API.] \code{grounded: true}, four citations, five context passages, and a trace with \code{total_ms: 37431.0}, the six stages, 2,367 input and 156 output tokens, \code{list_price_usd: 0.000299}, \code{retrieved: 20}, \code{cited: 4}. The response has no field that reveals the failed grading.
  \item[Demo.] The answer, four sources with their licence, a stage table totalling 37,431\,ms, \enquote{Cost at list price \$0.00030}, and under \emph{Notable}: \enquote{Self-check passed after 1 attempt(s)} \ev{app/app.py:132-139}.
  \item[CLI.] The answer, four \code{Sources} lines and the summary line beginning \code{37431ms [guard_input=0ms retrieve=1570ms rerank=14871ms guard_context=24ms generate=1782ms corrective=19184ms] 2367+156 tok}.
\end{description}

\section{Every function, in order}

\begin{longtable}{@{}>{\raggedright\arraybackslash}p{0.6cm}>{\raggedright\arraybackslash}p{5.9cm}>{\raggedright\arraybackslash}p{5.2cm}>{\raggedright\arraybackslash}p{3.6cm}@{}}
\toprule
\textbf{\#} & \textbf{Function} & \textbf{Where} & \textbf{Time / result}\\
\midrule
\endhead
1 & \code{Pipeline.answer} & \code{pipeline.py:144} & 19,184\,ms wall\\
2 & \code{QueryTrace(...)} & \code{pipeline.py:153} & ---\\
3 & \code{screen_input} $\rightarrow$ \code{scan} & \code{guard/input_guard.py:50}, \code{guard/patterns.py:133} & 0.1\,ms, clean\\
4 & \code{Pipeline._answer_corrective} & \code{pipeline.py:195} & ---\\
5 & \code{CorrectivePolicy(...)} & \code{correct/policy.py:46} & 0.8 / 0.5 / 2\\
6 & \code{run_corrective_loop} & \code{correct/loop.py:116} & 19,184\,ms (stage)\\
7 & \code{Pipeline.retrieve} $\rightarrow$ \code{retrieve_dense} $\rightarrow$ \code{embed_texts}, \code{query_points} & \code{pipeline.py:84}, \code{retrieve/dense.py:81}, \code{llm/provider.py:80} & 1,570\,ms, 20 passages\\
8 & \code{rerank} $\rightarrow$ \code{TextCrossEncoder.rerank}; \code{rank_movement} & \code{retrieve/rerank.py:44, 77} & 14,871\,ms\\
9 & \code{screen_context} & \code{guard/context_guard.py:61} & 24\,ms, all kept\\
10 & \code{generate_answer} $\rightarrow$ \code{build_answer_prompt} $\rightarrow$ \code{format_context} & \code{generate/answer.py:73}, \code{generate/prompts.py:58}, \code{generate/citations.py:54} & ---\\
11 & \code{client.models.generate_content} & Gemini \code{gemini-3.5-flash-lite} & 1,782\,ms, 2,367 + 156 tokens\\
12 & \code{_record_usage}, \code{_extract_text} & \code{generate/answer.py:41, 62} & ---\\
13 & \code{resolve_citations}, \code{strip_invalid_markers} & \code{generate/citations.py:81, 116} & markers 2--5 valid\\
14 & \code{grade_answer} $\rightarrow$ \code{generate_content} (JSON) & \code{correct/grader.py:180}, Gemini \code{gemini-3.1-flash-lite} & about 0.9\,s, 503 error\\
15 & \code{CorrectivePolicy.decide} & \code{correct/policy.py:54} & ACCEPT (rule 1)\\
16 & \code{LoopOutcome}, \code{Answer} & \code{correct/loop.py:162}, \code{pipeline.py:260} & returned\\
\bottomrule
\end{longtable}

\section{What this one trace demonstrates}

\begin{enumerate}
  \item \textbf{Reranking is the most expensive local stage}: 14.9 of 19.2 seconds here, and about 2 seconds in a steady-state benchmark --- either way far more than embedding the question (about 70\,ms) or searching the index (about 30\,ms).
  \item \textbf{The cross-encoder changed what the model read}: two of the five passages came from outside the dense top five, including a Glossary passage.
  \item \textbf{Citations were valid but not exhaustive}: the top-ranked passage was not cited, which the validator neither requires nor checks.
  \item \textbf{The double count is visible in a single query}: 37,431\,ms reported for 19,184\,ms of work.
  \item \textbf{Token accounting misses the grader}: 2,367 + 156 tokens recorded; the successful re-grade shows a grading call of 2,497 + 612.
  \item \textbf{Fail-open is real, and invisible}: a transient provider error turned the self-check off for this question, and every interface would have presented the answer exactly as if it had been verified.
  \item \textbf{Temperature 0 is not deterministic}: the same prompt produced 130 output tokens in one run and 156 in another.
\end{enumerate}
